% --- Archon physics formalization source begin ---
% archon:physics
% archon:covers PhyXMiniProblems/problem_phyx_mini_0939.lean
% archon:source-report reports/phyx_mini/problem_phyx_mini_0939.source.json

\chapter{Physics problem phyx\_mini\_0939}
\label{ch:PhyXMiniProblems_problem_phyx_mini_0939}

\paragraph{Problem source.}
\#\# Physical scenario

In one form of Tesla coil (a high-voltage generator popular in science museums), a long solenoid with length \textbackslash{}( l \textbackslash{}) and cross-sectional area \textbackslash{}( A \textbackslash{}) is closely wound with \textbackslash{}( N\_1 \textbackslash{}) turns of wire. A coil with \textbackslash{}( N\_2 \textbackslash{}) turns surrounds it at its center (figure).

\#\# Question

Find the mutual inductance \textbackslash{}( M \textbackslash{}).

\#\# Answer choices

- A: A: 22\textbackslash{} \textbackslash{}mu\textbackslash{}mathrm\{H\}

- B: B: 25\textbackslash{} \textbackslash{}mu\textbackslash{}mathrm\{H\}

- C: C: 75\textbackslash{} \textbackslash{}mu\textbackslash{}mathrm\{H\}

- D: D: 28\textbackslash{} \textbackslash{}mu\textbackslash{}mathrm\{H\}

\#\# Figure caption (auxiliary; use the image as primary evidence)

The image depicts a solenoid and an additional coil wrapped around it. The solenoid is shown as a series of closely spaced yellow wire coils, having a cross-sectional area labeled "A" and a length labeled "l." It has a certain number of turns indicated as \textbackslash{}( N\_1 \textbackslash{}).

Around this solenoid, there is another coil shown in gray, which has its turns labeled as \textbackslash{}( N\_2 \textbackslash{}). The gray coil wraps around the yellow solenoid, forming a secondary coil. The image includes lines to indicate and label the cross-sectional area, number of turns for each coil, and the length of the solenoid. The overall layout looks to explain a part of an electromagnetic or induction-related concept.

\#\# Dataset metadata

Electromagnetic Induction; Physical Model Grounding Reasoning, Multi-Formula Reasoning

\paragraph{Recorded answer/context.}
B: B: 25\textbackslash{} \textbackslash{}mu\textbackslash{}mathrm\{H\}

\paragraph{Figure/image path.}
/root/proposal\_for\_physic/science-mango/phyx\_mini\_run/phyx\_data/test\_image/939.png

\paragraph{Formalization target.}
create a compiling Lean file with sorry bodies at `PhyXMiniProblems/problem\_phyx\_mini\_0939.lean`. The Lean declarations must preserve the physical quantities, dimensions or dimensional roles, figure labels, governing-law hypotheses, and final relation expressed by this problem.
Use Mathlib/Physlib names found through LeanExplore where available. If a physics API is missing, introduce faithful local abstractions rather than scalar placeholder aliases.

\begin{theorem}[Physics formalization target]
\label{thm:physics:phyx_mini_0939:target}
\lean{PhyXMiniProblems.ProblemPhyXMini0939.problem_phyx_mini_0939}
\uses{def:physics:phyx-mini-0939:phyxminiproblems-problemphyxmini0939-coherentsireadout, def:physics:phyx-mini-0939:phyxminiproblems-problemphyxmini0939-teslacoilsetup, def:physics:phyx-mini-0939:phyxminiproblems-problemphyxmini0939-matcheswrittenteslacoilscenario, def:physics:phyx-mini-0939:phyxminiproblems-problemphyxmini0939-matchessuppliedteslacoilfigure, def:physics:phyx-mini-0939:phyxminiproblems-problemphyxmini0939-hasphysicalteslacoilparameters, def:physics:phyx-mini-0939:phyxminiproblems-problemphyxmini0939-usesfreespacepermeability, def:physics:phyx-mini-0939:phyxminiproblems-problemphyxmini0939-satisfiesidealteslacoillaws}
The assigned autoformalize agent should translate this physics problem into Lean declarations in the covered file, with theorem and lemma proof bodies written as `by sorry`.
\end{theorem}
\begin{proof}
This is an autoformalization task, not a proof task. Produce statements that can later be proved by the physics prover without weakening the physical model.
\end{proof}

% --- Archon declaration topology begin ---
\section{Declaration topology}
\begin{definition}[electric Current Dimension]
  \label{def:physics:phyx-mini-0939:phyxminiproblems-problemphyxmini0939-electriccurrentdimension}
  \lean{PhyXMiniProblems.ProblemPhyXMini0939.electricCurrentDimension}
  Electric current has dimension charge per time
\end{definition}
\begin{proof}
  This is the declared model definition of electric Current Dimension.
\end{proof}
\begin{definition}[area Dimension]
  \label{def:physics:phyx-mini-0939:phyxminiproblems-problemphyxmini0939-areadimension}
  \lean{PhyXMiniProblems.ProblemPhyXMini0939.areaDimension}
  Cross-sectional area has dimension length squared
\end{definition}
\begin{proof}
  This is the declared model definition of area Dimension.
\end{proof}
\begin{definition}[magnetic Flux Density Dimension]
  \label{def:physics:phyx-mini-0939:phyxminiproblems-problemphyxmini0939-magneticfluxdensitydimension}
  \lean{PhyXMiniProblems.ProblemPhyXMini0939.magneticFluxDensityDimension}
  Magnetic flux density has dimension M T⁻¹ C⁻¹
\end{definition}
\begin{proof}
  This is the declared model definition of magnetic Flux Density Dimension.
\end{proof}
\begin{definition}[magnetic Flux Dimension]
  \label{def:physics:phyx-mini-0939:phyxminiproblems-problemphyxmini0939-magneticfluxdimension}
  \lean{PhyXMiniProblems.ProblemPhyXMini0939.magneticFluxDimension}
  \uses{def:physics:phyx-mini-0939:phyxminiproblems-problemphyxmini0939-areadimension, def:physics:phyx-mini-0939:phyxminiproblems-problemphyxmini0939-magneticfluxdensitydimension}
  Magnetic flux and flux linkage have dimension B L²
\end{definition}
\begin{proof}
  This is the declared model definition of magnetic Flux Dimension.
\end{proof}
\begin{definition}[inductance Dimension]
  \label{def:physics:phyx-mini-0939:phyxminiproblems-problemphyxmini0939-inductancedimension}
  \lean{PhyXMiniProblems.ProblemPhyXMini0939.inductanceDimension}
  \uses{def:physics:phyx-mini-0939:phyxminiproblems-problemphyxmini0939-electriccurrentdimension, def:physics:phyx-mini-0939:phyxminiproblems-problemphyxmini0939-magneticfluxdimension}
  Inductance has the dimension of magnetic flux divided by current
\end{definition}
\begin{proof}
  This is the declared model definition of inductance Dimension.
\end{proof}
\begin{definition}[magnetic Permeability Dimension]
  \label{def:physics:phyx-mini-0939:phyxminiproblems-problemphyxmini0939-magneticpermeabilitydimension}
  \lean{PhyXMiniProblems.ProblemPhyXMini0939.magneticPermeabilityDimension}
  \uses{def:physics:phyx-mini-0939:phyxminiproblems-problemphyxmini0939-inductancedimension}
  Magnetic permeability has the dimension of inductance per length
\end{definition}
\begin{proof}
  This is the declared model definition of magnetic Permeability Dimension.
\end{proof}
\begin{definition}[Length Quantity]
  \label{def:physics:phyx-mini-0939:phyxminiproblems-problemphyxmini0939-lengthquantity}
  \lean{PhyXMiniProblems.ProblemPhyXMini0939.LengthQuantity}
  A nonnegative, unit-independent physical length
\end{definition}
\begin{proof}
  This is the declared model definition of Length Quantity.
\end{proof}
\begin{definition}[Area Quantity]
  \label{def:physics:phyx-mini-0939:phyxminiproblems-problemphyxmini0939-areaquantity}
  \lean{PhyXMiniProblems.ProblemPhyXMini0939.AreaQuantity}
  \uses{def:physics:phyx-mini-0939:phyxminiproblems-problemphyxmini0939-areadimension}
  A nonnegative, unit-independent physical area
\end{definition}
\begin{proof}
  This is the declared model definition of Area Quantity.
\end{proof}
\begin{definition}[Electric Current Quantity]
  \label{def:physics:phyx-mini-0939:phyxminiproblems-problemphyxmini0939-electriccurrentquantity}
  \lean{PhyXMiniProblems.ProblemPhyXMini0939.ElectricCurrentQuantity}
  \uses{def:physics:phyx-mini-0939:phyxminiproblems-problemphyxmini0939-electriccurrentdimension}
  A nonnegative, unit-independent electric-current magnitude
\end{definition}
\begin{proof}
  This is the declared model definition of Electric Current Quantity.
\end{proof}
\begin{definition}[Magnetic Flux Density Quantity]
  \label{def:physics:phyx-mini-0939:phyxminiproblems-problemphyxmini0939-magneticfluxdensityquantity}
  \lean{PhyXMiniProblems.ProblemPhyXMini0939.MagneticFluxDensityQuantity}
  \uses{def:physics:phyx-mini-0939:phyxminiproblems-problemphyxmini0939-magneticfluxdensitydimension}
  A nonnegative, unit-independent magnetic-flux-density magnitude
\end{definition}
\begin{proof}
  This is the declared model definition of Magnetic Flux Density Quantity.
\end{proof}
\begin{definition}[Magnetic Flux Quantity]
  \label{def:physics:phyx-mini-0939:phyxminiproblems-problemphyxmini0939-magneticfluxquantity}
  \lean{PhyXMiniProblems.ProblemPhyXMini0939.MagneticFluxQuantity}
  \uses{def:physics:phyx-mini-0939:phyxminiproblems-problemphyxmini0939-magneticfluxdimension}
  A nonnegative magnetic flux or flux linkage
\end{definition}
\begin{proof}
  This is the declared model definition of Magnetic Flux Quantity.
\end{proof}
\begin{definition}[Magnetic Permeability Quantity]
  \label{def:physics:phyx-mini-0939:phyxminiproblems-problemphyxmini0939-magneticpermeabilityquantity}
  \lean{PhyXMiniProblems.ProblemPhyXMini0939.MagneticPermeabilityQuantity}
  \uses{def:physics:phyx-mini-0939:phyxminiproblems-problemphyxmini0939-magneticpermeabilitydimension}
  A nonnegative, unit-independent magnetic permeability
\end{definition}
\begin{proof}
  This is the declared model definition of Magnetic Permeability Quantity.
\end{proof}
\begin{definition}[Mutual Inductance Quantity]
  \label{def:physics:phyx-mini-0939:phyxminiproblems-problemphyxmini0939-mutualinductancequantity}
  \lean{PhyXMiniProblems.ProblemPhyXMini0939.MutualInductanceQuantity}
  \uses{def:physics:phyx-mini-0939:phyxminiproblems-problemphyxmini0939-inductancedimension}
  A nonnegative, unit-independent mutual inductance
\end{definition}
\begin{proof}
  This is the declared model definition of Mutual Inductance Quantity.
\end{proof}
\begin{definition}[coherent SIReadout]
  \label{def:physics:phyx-mini-0939:phyxminiproblems-problemphyxmini0939-coherentsireadout}
  \lean{PhyXMiniProblems.ProblemPhyXMini0939.coherentSIReadout}
  Read a nonnegative dimensionful quantity in Physlib's coherent SI units
\end{definition}
\begin{proof}
  This is the declared model definition of coherent SIReadout.
\end{proof}
\begin{definition}[Winding Role]
  \label{def:physics:phyx-mini-0939:phyxminiproblems-problemphyxmini0939-windingrole}
  \lean{PhyXMiniProblems.ProblemPhyXMini0939.WindingRole}
  The two windings distinguished by the prose and primary raster
\end{definition}
\begin{proof}
  This is the declared model definition of Winding Role.
\end{proof}
\begin{definition}[Figure Label]
  \label{def:physics:phyx-mini-0939:phyxminiproblems-problemphyxmini0939-figurelabel}
  \lean{PhyXMiniProblems.ProblemPhyXMini0939.FigureLabel}
  The four symbolic annotations visible in 939.png
\end{definition}
\begin{proof}
  This is the declared model definition of Figure Label.
\end{proof}
\begin{definition}[Winding Color]
  \label{def:physics:phyx-mini-0939:phyxminiproblems-problemphyxmini0939-windingcolor}
  \lean{PhyXMiniProblems.ProblemPhyXMini0939.WindingColor}
  The winding colors used in the primary raster
\end{definition}
\begin{proof}
  This is the declared model definition of Winding Color.
\end{proof}
\begin{definition}[Primary Winding Model]
  \label{def:physics:phyx-mini-0939:phyxminiproblems-problemphyxmini0939-primarywindingmodel}
  \lean{PhyXMiniProblems.ProblemPhyXMini0939.PrimaryWindingModel}
  Physical idealizations available for the primary winding
\end{definition}
\begin{proof}
  This is the declared model definition of Primary Winding Model.
\end{proof}
\begin{definition}[Secondary Placement]
  \label{def:physics:phyx-mini-0939:phyxminiproblems-problemphyxmini0939-secondaryplacement}
  \lean{PhyXMiniProblems.ProblemPhyXMini0939.SecondaryPlacement}
  Placements available for the secondary relative to the primary
\end{definition}
\begin{proof}
  This is the declared model definition of Secondary Placement.
\end{proof}
\begin{definition}[Tesla Coil Figure]
  \label{def:physics:phyx-mini-0939:phyxminiproblems-problemphyxmini0939-teslacoilfigure}
  \lean{PhyXMiniProblems.ProblemPhyXMini0939.TeslaCoilFigure}
  \uses{def:physics:phyx-mini-0939:phyxminiproblems-problemphyxmini0939-windingrole, def:physics:phyx-mini-0939:phyxminiproblems-problemphyxmini0939-figurelabel, def:physics:phyx-mini-0939:phyxminiproblems-problemphyxmini0939-windingcolor}
  Literal presentation data from the supplied raster, with no numeric data
\end{definition}
\begin{proof}
  This is the declared model definition of Tesla Coil Figure.
\end{proof}
\begin{definition}[Tesla Coil Setup]
  \label{def:physics:phyx-mini-0939:phyxminiproblems-problemphyxmini0939-teslacoilsetup}
  \lean{PhyXMiniProblems.ProblemPhyXMini0939.TeslaCoilSetup}
  \uses{def:physics:phyx-mini-0939:phyxminiproblems-problemphyxmini0939-lengthquantity, def:physics:phyx-mini-0939:phyxminiproblems-problemphyxmini0939-areaquantity, def:physics:phyx-mini-0939:phyxminiproblems-problemphyxmini0939-electriccurrentquantity, def:physics:phyx-mini-0939:phyxminiproblems-problemphyxmini0939-magneticfluxdensityquantity, def:physics:phyx-mini-0939:phyxminiproblems-problemphyxmini0939-magneticfluxquantity, def:physics:phyx-mini-0939:phyxminiproblems-problemphyxmini0939-magneticpermeabilityquantity, def:physics:phyx-mini-0939:phyxminiproblems-problemphyxmini0939-mutualinductancequantity, def:physics:phyx-mini-0939:phyxminiproblems-problemphyxmini0939-primarywindingmodel, def:physics:phyx-mini-0939:phyxminiproblems-problemphyxmini0939-secondaryplacement, def:physics:phyx-mini-0939:phyxminiproblems-problemphyxmini0939-teslacoilfigure}
  Independent physical data for the idealized experiment. In particular, mutualInductance is not defined from the requested closed form; it is related to the other observables only by the governing-law premise below
\end{definition}
\begin{proof}
  This is the declared model definition of Tesla Coil Setup.
\end{proof}
\begin{definition}[Matches Written Tesla Coil Scenario]
  \label{def:physics:phyx-mini-0939:phyxminiproblems-problemphyxmini0939-matcheswrittenteslacoilscenario}
  \lean{PhyXMiniProblems.ProblemPhyXMini0939.MatchesWrittenTeslaCoilScenario}
  \uses{def:physics:phyx-mini-0939:phyxminiproblems-problemphyxmini0939-teslacoilsetup}
  The qualitative winding arrangement stated in the problem prose
\end{definition}
\begin{proof}
  This is the declared model definition of Matches Written Tesla Coil Scenario.
\end{proof}
\begin{definition}[Matches Supplied Tesla Coil Figure]
  \label{def:physics:phyx-mini-0939:phyxminiproblems-problemphyxmini0939-matchessuppliedteslacoilfigure}
  \lean{PhyXMiniProblems.ProblemPhyXMini0939.MatchesSuppliedTeslaCoilFigure}
  \uses{def:physics:phyx-mini-0939:phyxminiproblems-problemphyxmini0939-teslacoilsetup}
  Literal label, color, and placement evidence from the primary raster
\end{definition}
\begin{proof}
  This is the declared model definition of Matches Supplied Tesla Coil Figure.
\end{proof}
\begin{definition}[Has Physical Tesla Coil Parameters]
  \label{def:physics:phyx-mini-0939:phyxminiproblems-problemphyxmini0939-hasphysicalteslacoilparameters}
  \lean{PhyXMiniProblems.ProblemPhyXMini0939.HasPhysicalTeslaCoilParameters}
  \uses{def:physics:phyx-mini-0939:phyxminiproblems-problemphyxmini0939-coherentsireadout, def:physics:phyx-mini-0939:phyxminiproblems-problemphyxmini0939-teslacoilsetup}
  Positivity and nondegeneracy of the physical parameters
\end{definition}
\begin{proof}
  This is the declared model definition of Has Physical Tesla Coil Parameters.
\end{proof}
\begin{definition}[Uses Free Space Permeability]
  \label{def:physics:phyx-mini-0939:phyxminiproblems-problemphyxmini0939-usesfreespacepermeability}
  \lean{PhyXMiniProblems.ProblemPhyXMini0939.UsesFreeSpacePermeability}
  \uses{def:physics:phyx-mini-0939:phyxminiproblems-problemphyxmini0939-coherentsireadout, def:physics:phyx-mini-0939:phyxminiproblems-problemphyxmini0939-teslacoilsetup}
  The hollow air-core idealization identifies the dimensionful permeability's coherent-SI readout with Physlib's free-space system parameter μ₀. This premise contains no mutual-inductance formula or answer value
\end{definition}
\begin{proof}
  This is the declared model definition of Uses Free Space Permeability.
\end{proof}
\begin{definition}[Satisfies Ideal Tesla Coil Laws]
  \label{def:physics:phyx-mini-0939:phyxminiproblems-problemphyxmini0939-satisfiesidealteslacoillaws}
  \lean{PhyXMiniProblems.ProblemPhyXMini0939.SatisfiesIdealTeslaCoilLaws}
  \uses{def:physics:phyx-mini-0939:phyxminiproblems-problemphyxmini0939-coherentsireadout, def:physics:phyx-mini-0939:phyxminiproblems-problemphyxmini0939-teslacoilsetup}
  The four independent textbook laws used in the derivation. They are stated at the coherent-SI readout boundary and do not contain the requested closed form M = μ₀ N₁ N₂ A / l
\end{definition}
\begin{proof}
  This is the declared model definition of Satisfies Ideal Tesla Coil Laws.
\end{proof}
\begin{definition}[Answer Choice]
  \label{def:physics:phyx-mini-0939:phyxminiproblems-problemphyxmini0939-answerchoice}
  \lean{PhyXMiniProblems.ProblemPhyXMini0939.AnswerChoice}
  The four answer labels printed with the problem
\end{definition}
\begin{proof}
  This is the declared model definition of Answer Choice.
\end{proof}
\begin{definition}[displayed Microhenries]
  \label{def:physics:phyx-mini-0939:phyxminiproblems-problemphyxmini0939-answerchoice-displayedmicrohenries}
  \lean{PhyXMiniProblems.ProblemPhyXMini0939.AnswerChoice.displayedMicrohenries}
  \uses{def:physics:phyx-mini-0939:phyxminiproblems-problemphyxmini0939-answerchoice}
  Each answer choice's displayed mutual-inductance value, in microhenries
\end{definition}
\begin{proof}
  This is the declared model definition of displayed Microhenries.
\end{proof}
\begin{definition}[recorded Dataset Answer]
  \label{def:physics:phyx-mini-0939:phyxminiproblems-problemphyxmini0939-recordeddatasetanswer}
  \lean{PhyXMiniProblems.ProblemPhyXMini0939.recordedDatasetAnswer}
  \uses{def:physics:phyx-mini-0939:phyxminiproblems-problemphyxmini0939-answerchoice}
  The source dataset records answer label B; this is metadata, not a law
\end{definition}
\begin{proof}
  This is the declared model definition of recorded Dataset Answer.
\end{proof}
% --- Archon declaration topology end ---
% --- Archon physics formalization source end ---
